\documentclass[14pt]{report}

\usepackage[utf8]{inputenc} % Enables LaTeX to handle Unicode Characters directly 
\usepackage{mathptmx} % Used to change the default font in math mode
\usepackage{amsmath} % To write mathematical equation
\usepackage{graphicx} % Used to include external graphics
\usepackage{animate} % for animation
\usepackage{subcaption} % To present multiple related images or tables within a single figure
\usepackage{booktabs} % to creat high quality tables
\usepackage{circuitikz} % to draw electrical circuit
\usepackage{siunitx} % for si unit
\usepackage{tikz} % to draaw flow chart
\usetikzlibrary{shapes.geometric, arrows} % for flow chart
\usepackage{float} % enhanced control over the placement of floating element
\usepackage{enumerate} % to customize the appereance of lists
\usepackage[colorlinks,linkcolor=black, citecolor=blue, urlcolor=blue]{hyperref} % to enable the creation of hyperlinks
\usepackage{geometry} % for margin specification
\geometry{left=3.5cm,right=2.5cm,top=2.5cm,bottom=2.5cm} % set margin
\usepackage{setspace} % set space between line
\onehalfspacing
% \usepackage{fontspec} % to utilize fonts installed on system directly
% \setmainfont{Times New Roman}
\usepackage{titlesec}
\usepackage{fancyhdr} % to customize header and footer
\usepackage{lipsum} % to generate dummy text
\usepackage{lastpage} % page numbering
\usepackage{etoolbox} % provides tools for modifying and patching LaTeX commands. 
\usepackage{cite} % of citations and bibliographies in LaTeX.

\usepackage{lmodern} % to control font-size
%% {\fontsize{font size}{base line strech} \selectfont}
%% {\fontsize{40}{48} \selectfont Dynamic} % syntex to control fontsize.


%%%%%%%%%%%%%%%%%%%%%%% Coding Style %%%%%%%%%%%%%%%%%%%%%%%%%
\usepackage{listings}
\usepackage{xcolor}

%New colors defined below
\definecolor{codegreen}{rgb}{0,0.6,0}
\definecolor{codegray}{rgb}{0.5,0.5,0.5}
\definecolor{codepurple}{rgb}{0.58,0,0.82}
\definecolor{backcolour}{rgb}{0.95,0.95,0.92}

%Code listing style named "mystyle"
\lstdefinestyle{mystyle}{
  backgroundcolor=\color{backcolour}, commentstyle=\color{codegreen},
  keywordstyle=\color{magenta},
  numberstyle=\tiny\color{codegray},
  stringstyle=\color{codepurple},
  basicstyle=\ttfamily\small,
  breakatwhitespace=false,         
  breaklines=true,                 
  captionpos=b,                    
  keepspaces=true,                 
  numbers=left,                    
  numbersep=5pt,                  
  showspaces=false,                
  showstringspaces=false,
  showtabs=false,                  
  tabsize=2
}

\lstset{style=mystyle}

%%%%%%%%%%%%%%%%%%%%%%%%%%%%%%%%%%%
\usepackage{tabularx}
\usepackage{makecell}
\usepackage{multirow}
\usepackage{multicol}
\usepackage{hhline}
\usepackage{pgf-pie}
\usetikzlibrary{backgrounds}
\usepackage{courier}
\newcommand\tab[1][1cm]{\hspace*{#1}}

%%%%%% To insert PDF cover page %%%%%%%
\usepackage{titlepic}
\usepackage{pdfpages}

\title{My Knowledge On Java}
\author{Aong Cho Marma}
\date{13 July, 2023}

\begin{document}
	\includepdf[pages=1, fitpaper]{CoverPage/CoverPage}
	
	\maketitle
	
	%%%%%%%%%%%%%%%%%%% Table of Contents %%%%%%%%%%%%%%%%
	\tableofcontents
	\pagenumbering{\roman{i}}	
	
	%%%%%%%%%%%%%%%%%% Page Numbering Style %%%%%%%%%%%%%%
	\pagenumbering{arabic}
	\patchcmd{\chapter}{\thispagestyle{plain}}{\thispagestyle{fancy}}{}{}
	\fancypagestyle{main}{
    	\fancyhf{} % clear all header and footer fields
    	\renewcommand{\headrulewidth}{0pt} % remove header rule
    	\renewcommand{\footrulewidth}{1pt} % set footerr rule
    	\fancyfoot[L]{
    		\href{mailto:aongcho880@gmail.com}{\underline{Aong Cho Marma}}\\
    		\fontsize{11}{13}\selectfont \leftmark
    	} % chapter name at left footer
    	\fancyfoot[R]{Page \thepage\ of \pageref{LastPage}} % page number at center footer
	}
	
	\pagestyle{main} % set the main page style
	\setcounter{page}{1}
	% Define style for equation numbers
	\renewcommand{\theequation}{\arabic{chapter}.\arabic{equation}}
	
	%%%%%%%%%%%%%%%%%%%%%%%%%%%%%%%%%%%%%%%%%%%%%
	
	%%%%%%%%%%%%%%% Chapters %%%%%%%%%%%%%%%%%%%%
	%\newpage
\chapter{INTRODUCTION}
\section{Basic Competitive Programming}
Competitive programming is a sport of coding where participants solve complex algorithmic problems under time constraints. Using programming languages, they devise efficient solutions that navigate data structures, math, and logic. This intellectually stimulating activity hones problem-solving skills and fosters creativity, often practiced through online platforms and contests.
	\chapter{File Handling In Java}

% ---------------- Section 1 -----------------------
\section{Create, Delete and Get The Path Of The Dir.}
\lstset{
  basicstyle=\fontsize{8}{10}\selectfont\ttfamily,
}
\begin{lstlisting}[language=java]
package fileHandling;

import java.io.File;

public class A_CreateDir {

	public static void main(String[] args) {
		// It will create the directory at the current project directory.
		File dir = new File("Test"); // we can use directory path also
		dir.mkdir();
		
		//---- Get The File Directory
		String dirPath = dir.getAbsolutePath();
		String name = dir.getName();
		System.out.println("Dir Name: "+name+"\nPath: "+dirPath);
		
		//---- Delete Directory If Exist
		if(dir.delete()) {
			System.out.println(name+" directory has been deleted.");
		}
		
	}
}

\end{lstlisting}


% ---------------- Section 2 -----------------------
\newpage
\section{Create File and Delete File}
\lstset{
  	basicstyle=\fontsize{8}{10}\selectfont\ttfamily,
}
\begin{lstlisting}[language=java]
package fileHandling;

import java.io.File;
import java.util.*;

public class B_CreateFile {

	public static void main(String[] args) {
		//---- File Must Be Exist Other Wise Create Directory
		//---- Set The Directory path
		File dir = new File("TextFile");
		
		//------ Get The Directory Path
		String path = dir.getAbsolutePath();
		
		//---- Set File Name And Path
		File file1 = new File(path+"/file1.txt");
		File file2 = new File(path+"/file2.txt");
		
		try {
			
			// ---- Create New File
			file1.createNewFile();
			file2.createNewFile();
			
		} catch (Exception e) {
			System.out.println(e);
		}
		
		//--- Check If A File Exist
		if(file2.exists()) {
			//---- Delete File
			file2.delete();
		}
	}
}

\end{lstlisting}


%------------- Ways to write data into a file -----------------
\newpage
\section{Ways To Write Data Into a File in Java}

\subsection{FileWriter Class}
\begin{itemize}
	\item \textit{FileWriter} is used for writing character data into a file.
	\item It`s suitable for writing simple text-based data.
	\item You can write strings and characters directly to the file.
\end{itemize}
\textbf{Code:}
\lstset{
  	basicstyle=\fontsize{8}{10}\selectfont\ttfamily,
}
\begin{lstlisting}[language=java]
package fileHandling;

import java.io.File;
import java.io.FileWriter;
import java.io.IOException;
import java.util.Formatter;

public class WriteInFile {
	public static void main(String[] args) {	
		FileWriter writer = null;
		
		try {
			writer = new FileWriter("Output.txt");
			
			char c = 'A';
			writer.write(c); //----- write single character
			writer.write('\n');
			
			char[] charArry = {'X', 'Y', 'Z'};
			writer.write(charArry); //--- Write array of characters
			writer.write('\n');
			
			String str = "This is a string";
			writer.write(str);	//--- write string
			writer.write('\n');
			
			writer.flush(); //- Ensure the data is immediately written to the file
			writer.close();
		} catch (IOException e) {
			System.out.println(e);
		}
	}
}
\end{lstlisting}

%--------------------------------
\newpage
\subsection{BufferedWriter Class}
\begin{itemize}
	\item This approach combines a \textbf{BufferedWriter} with a \textbf{FileWriter} to improve efficiency when writing large amounts of text data.
	\item It reduces the number of disk writes and is useful for optimizing performance
\end{itemize}
\textbf{Code:}
\lstset{
  	basicstyle=\fontsize{8}{10}\selectfont\ttfamily,
}
\begin{lstlisting}[language=java]
package fileHandling;

import java.io.BufferedWriter;
import java.io.FileWriter;
import java.io.IOException;

public class WriteInFile {

	public static void main(String[] args) {
		
		FileWriter fWriter = null;
		BufferedWriter bWriter = null;
		
		try {
			fWriter = new FileWriter("Output.txt");
			bWriter = new BufferedWriter(fWriter);
			
			bWriter.write('A');
			bWriter.newLine();
			bWriter.write("This is a String");
			bWriter.flush();
			
			fWriter.close();
			bWriter.close();
		} catch (IOException e) {
			System.out.println(e);
		}
	}

}

\end{lstlisting}


%--------------------------------
\newpage
\subsection{PrintWriter Class}
\begin{itemize}
	\item \textbf{PrintWriter} is useful for writing formatted text data into a file.
	\item It provides methods like \textbf{\texttt{printf}} and \textbf{\texttt{println}} for formatting output.
\end{itemize}

\textbf{Code:}
\lstset{
  	basicstyle=\fontsize{8}{10}\selectfont\ttfamily,
}
\begin{lstlisting}[language=java]
package fileHandling;

import java.io.FileWriter;
import java.io.IOException;
import java.io.PrintWriter;

public class WriteInFile {

	public static void main(String[] args) {
		
		FileWriter fWriter = null;
		PrintWriter pWriter = null;
		
		try {
			fWriter = new FileWriter("Output.txt");
			pWriter = new PrintWriter(fWriter);
			

			pWriter.println("Hello World");
			pWriter.println("This is PrintWriter");
			pWriter.printf("Formatted Output: %s : %d + %d = %d %n", "Sum",5,7,12);
			pWriter.write("New Line");
			pWriter.flush();
			
			fWriter.close();
			pWriter.close();
		} catch (IOException e) {
			System.out.println(e);
		}
	}

}

\end{lstlisting}


%-------------- Java Formatter Class ----------------
\newpage
\section{Java Formatter Class}
The Java \textit{Formatter} class is defined in the \textit{java.util} package and is declared final. It, therefore, cannot be extended or sub-classed.\\
With the help of this class, we can send formatted outputs to other outputs streams or devices, such as a GUI component or to a file apart from standard output.\\
\textbf{Formatter Construction}
\begin{itemize}
	\item \textbf{Formatter() :}
	\item \textbf{Formatter(Appendable a) :}
	\item \textbf{Formatter(Appendable a, Locale loc) :}
	\item \textbf{Formatter(File file) :} The file parameter of this constructor designates a reference to a open file where the output will be streamed.\\
\end{itemize}

\textbf{Using Formatter}
\begin{itemize}
	\item \textbf{\%S or \%s :} Specifies String
	\item \textbf{\%X or \%x :} Specifies hexadecimal integer
	\item \textbf{\%o :} Specifies Octal integer
	\item \textbf{\%d :} Specifies Decimal integer
	\item \textbf{\%c :} Specifies character
	\item \textbf{\%T or \%t :} Specifies Time and Date
	\item \textbf{\%n :} Insets newline character
	\item \textbf{\%B or \%b :} Specifies Boolean
	\item \textbf{\%A or \%a :} Specifies floating point hexadecimal
	\item \textbf{\%f :} Specifies Decimal floating point\\
\end{itemize}

\newpage
\subsection{A Few Quick Examples}

\textbf{Using argument\_index}

\lstset{
  	basicstyle=\fontsize{8}{10}\selectfont\ttfamily,
}
\begin{lstlisting}[language=java]
Formatter f2 = new Formatter();
f2.format("%2$s %1$s %3$s", "fear", "weakness","strengthen");
System.out.println(f2);
f2.close();

// Output: weakness fear strengthen		
\end{lstlisting}


\textbf{Regionalize Date}

\lstset{
  	basicstyle=\fontsize{8}{10}\selectfont\ttfamily,
}
\begin{lstlisting}[language=java]
Formatter f3=new Formatter();
f3.format(Locale.FRENCH,"%1$te %1$tB, %1$tY", Calendar.getInstance());
System.out.println(f3);
f3.close();

Formatter f4=new Formatter();
f4.format(Locale.ENGLISH,"%1$te %1$tB, %1$tY",Calendar.getInstance());
System.out.println(f4);
f4.close();

// Output: 7 octobre, 2023
//  		   7 October, 2023		
\end{lstlisting}


\textbf{Using \%n and \%\% Specifiers}

\lstset{
  	basicstyle=\fontsize{8}{10}\selectfont\ttfamily,
}
\begin{lstlisting}[language=java]
Formatter f = new Formatter();
f.format("Format%n %.2f%% complete", 46.6);
System.out.println(f);
f.close();
	
// Output: Format
//			46.60% complete
\end{lstlisting}

\textbf{Write in a file}

\lstset{
  	basicstyle=\fontsize{8}{10}\selectfont\ttfamily,
}
\begin{lstlisting}[language=java]
public class WriteInFile {
	public static void main(String[] args) {
		//--- "TestFile" directory must be exist in the project directory
		File file = new File("TextFile");
		String path = file.getAbsolutePath();
		System.out.println(path);
		
		try {
			Formatter formatter = new Formatter(path+"/file1.txt");
			formatter.format("%s %s %s\n", "21701002","Aong Cho","CSE");
			formatter.format("%s %s %s\n", "21701001","Taqi Ismile","CSE");
			formatter.close();
		} catch (Exception e) {
			System.out.println(e);
		}
	}
}
\end{lstlisting}






	
	
\end{document}